% =====================================================================
%  Chapitre 2 — Conception
% =====================================================================
\chapter{Conception du système}
\label{chap:conception}

\section{Acteurs et rôles}

Le modèle de droits distingue onze rôles fonctionnels. Chacun combine un jeu de
permissions et un niveau de périmètre, c'est-à-dire l'étendue des données sur
lesquelles il agit. Les descriptions qui suivent reprennent le catalogue défini dans
\code{app/core/rbac.py}.

\subsection{Administrateur Medicacom}

\textbf{Mission.} Administrer la plateforme et les organisations clientes : référentiels,
modèles de questionnaires, rôles disponibles et support des tenants.

\textbf{Périmètre.} Plateforme — le seul rôle qui franchit la frontière des
organisations.

\textbf{Droits.} Les 41~permissions du catalogue, dont \code{platform.admin}, réservée à
ce rôle. Un administrateur d'entreprise ne peut ni se l'attribuer ni la placer dans un
rôle qu'il crée.

\subsection{Administrateur entreprise}

\textbf{Mission.} Administrer son organisation : sites, départements, utilisateurs,
rôles et paramètres.

\textbf{Périmètre.} Organisation. Il n'accède jamais à une autre.

\textbf{Droits.} 23~permissions : toutes les lectures, plus la gestion de
l'organisation, des sites, des départements, des comptes et des rôles, ainsi que
l'accès au journal d'audit.

\textbf{Limitations.} Il n'arbitre pas le portefeuille et ne renseigne pas les
diagnostics. Son écran d'accueil est la console d'administration, pas un tableau de
bord.

\subsection{Direction Générale}

\textbf{Mission.} Vision globale du portefeuille, arbitrage et décision de comité.

\textbf{Périmètre.} Organisation.

\textbf{Droits.} 22~permissions : toutes les lectures, le tableau de bord Direction, la
validation d'opportunité, la repriorisation du portefeuille, le passage de gate,
la décision de comité sur la fiche projet et la validation des bénéfices financiers.

\textbf{Limitations.} La Direction n'administre pas les comptes et ne renseigne pas le
questionnaire de diagnostic. Elle le consulte.

\subsection{Responsable de département}

\textbf{Mission.} Piloter les besoins et les projets de son ou ses départements.

\textbf{Périmètre.} Département — le seul rôle dont la portée est restreinte par
rattachement.

\textbf{Droits.} 18~permissions : création et modification d'opportunités, réponse et
validation du diagnostic, rédaction de fiche projet, tableau de bord de département.

\textbf{Limitations.} Il consulte les dépendances transverses mais ne modifie rien hors
de ses départements. Une tentative d'écriture ailleurs répond 403.

\subsection{Contributeur métier}

\textbf{Mission.} Répondre aux questionnaires assignés, documenter les processus et
proposer des opportunités.

\textbf{Périmètre.} Contributions personnelles.

\textbf{Droits.} 11~permissions, le jeu le plus restreint : lecture de l'organisation et
de son département, réponse au diagnostic, création et lecture d'opportunités, lecture
des recommandations et des fiches.

\textbf{Limitations.} Aucun accès au portefeuille consolidé ni aux données financières.
Aucun arbitrage, aucune validation.

\subsection{DSI / Architecte SI}

\textbf{Mission.} Évaluer la faisabilité technique, tenir l'architecture AS-IS et TO-BE
ainsi que les dépendances applicatives.

\textbf{Droits.} 16~permissions, dont \code{architecture.manage} qui lui est propre avec
l'administrateur, et \code{project.validate\_gate} pour les gates techniques.

\subsection{Data Office / Data Engineer}

\textbf{Mission.} Tenir le référentiel des sources de données, leur qualité, leur
ownership et les dépendances data interservices.

\textbf{Droits.} 14~permissions, dont \code{data.manage}. C'est le seul rôle, avec
l'administrateur de plateforme, qui écrit dans le référentiel des sources.

\textbf{Limitations.} Il ne gère pas le référentiel applicatif, qui relève du DSI.

\subsection{AI Lead / Data Scientist}

\textbf{Mission.} Analyser les opportunités qualifiées, amender la recommandation avec
justification, définir les critères de performance et les garde-fous.

\textbf{Droits.} 16~permissions, dont \code{recommendation.override}, qui lui est
exclusive. C'est le seul rôle autorisé à relancer la qualification d'une opportunité.

\textbf{Limitations.} Il n'administre pas les comptes.

\subsection{Finance / Contrôle de gestion}

\textbf{Mission.} Instruire CAPEX, OPEX et gains, valider les hypothèses du business
case, comparer valeur attendue et valeur réalisée.

\textbf{Droits.} 15~permissions, dont \code{finance.manage} et \code{finance.validate}.

\textbf{Limitations.} Finance ne touche ni à la recommandation IA ni à l'architecture
technique. Un appel à l'endpoint de qualification depuis un compte Finance répond 403.

\subsection{PMO / Chef de projet}

\textbf{Mission.} Gérer le cycle de vie des projets : phases, jalons, tâches, risques,
incidents, budget de suivi et reporting.

\textbf{Droits.} 15~permissions, dont \code{project.manage}.

\textbf{Limitations.} Il ne valide pas les bénéfices financiers.

\subsection{Auditeur / Lecture seule}

\textbf{Mission.} Consulter preuves, scores, versions, décisions et journal d'audit.

\textbf{Droits.} 16~permissions, toutes en lecture, plus \code{audit.read}.

\textbf{Limitations.} Aucune création, modification, suppression ni validation, quel que
soit le module. Huit tests paramétrés vérifient ce point endpoint par endpoint.

\subsection{Matrice de synthèse}

\begin{table}[H]
\centering
\caption{Capacités par rôle, dérivées du catalogue de permissions}
\label{tab:matrice-roles}
\small
\begin{tabular}{@{}L{4.4cm} C{1.6cm} C{1.5cm} C{1.8cm} C{1.7cm} C{1.9cm}@{}}
\toprule
\textbf{Rôle} & \textbf{Consul\-tation} & \textbf{Créa\-tion} &
\textbf{Modifi\-cation} & \textbf{Valida\-tion} & \textbf{Adminis\-tration} \\
\midrule
Administrateur Medicacom   & \oui & \oui & \oui & \oui & \oui \\
Administrateur entreprise  & \oui & \oui & \oui & \non & \oui \\
Direction Générale         & \oui & \non & \oui & \oui & \non \\
Responsable de département & \oui & \oui & \oui & \oui & \non \\
Contributeur métier        & \oui & \oui & \oui & \non & \non \\
DSI / Architecte SI        & \oui & \oui & \oui & \oui & \non \\
Data Office / Data Engineer& \oui & \oui & \oui & \non & \non \\
AI Lead / Data Scientist   & \oui & \oui & \oui & \oui & \non \\
Finance / Contrôle gestion & \oui & \oui & \oui & \oui & \non \\
PMO / Chef de projet       & \oui & \oui & \oui & \non & \non \\
Auditeur / Lecture seule   & \oui & \non & \non & \non & \non \\
\bottomrule
\end{tabular}
\end{table}

La colonne \emph{Validation} recouvre des objets différents selon le rôle, ce que le
tableau ne peut pas montrer. La Direction valide les opportunités, les gates, les
fiches et les bénéfices. Le responsable de département ne valide que les campagnes de
diagnostic. Le DSI et l'AI Lead ne valident que les gates de projet. Finance ne valide
que les bénéfices financiers. La granularité réelle est celle de la permission, pas
celle de la colonne.

\section{Cas d'utilisation}

\subsection{Vue générale}

\begin{figure}[H]
\centering
\begin{tikzpicture}[scale=0.92, transform shape]
  % --- acteurs à gauche
  \node[acteur] (dir)   at (0,4.6)  {Direction Générale};
  \node[acteur] (resp)  at (0,3.0)  {Responsable de département};
  \node[acteur] (contrib) at (0,1.4) {Contributeur métier};
  \node[acteur] (adm)   at (0,-0.2) {Administrateur entreprise};

  % --- acteurs à droite
  \node[acteur] (dsi)   at (12,4.6) {DSI / Architecte};
  \node[acteur] (data)  at (12,3.0) {Data Office};
  \node[acteur] (ai)    at (12,1.4) {AI Lead};
  \node[acteur] (fin)   at (12,-0.2){Finance};
  \node[acteur] (pmo)   at (12,-1.8){PMO};
  \node[acteur] (audit) at (0,-1.8) {Auditeur};

  % --- système
  \begin{scope}[on background layer]
    \draw[draw=medicaBorder, fill=medicaBg, rounded corners=8pt]
      (2.7,-2.9) rectangle (9.3,5.6);
  \end{scope}
  \node[font=\small\bfseries, text=medicaDeep] at (6,5.25) {Medicacom AOI Transformation};

  \node[cas] (connect) at (6,4.45) {Se connecter};
  \node[cas] (diag)    at (4.4,3.5) {Réaliser le diagnostic};
  \node[cas] (repond)  at (7.6,3.5) {Répondre au questionnaire};
  \node[cas] (opp)     at (4.4,2.5) {Créer une opportunité};
  \node[cas] (qualif)  at (7.6,2.5) {Qualifier l'opportunité};
  \node[cas] (reco)    at (4.4,1.5) {Consulter la recommandation};
  \node[cas] (over)    at (7.6,1.5) {Amender la recommandation};
  \node[cas] (prio)    at (4.4,0.5) {Prioriser le portefeuille};
  \node[cas] (fiche)   at (7.6,0.5) {Générer la fiche projet};
  \node[cas] (valid)   at (4.4,-0.5) {Valider en comité};
  \node[cas] (suivi)   at (7.6,-0.5) {Suivre le projet};
  \node[cas] (roi)     at (4.4,-1.5) {Consulter le ROI};
  \node[cas] (admin)   at (7.6,-1.5) {Administrer l'organisation};
  \node[cas] (journal) at (6,-2.4) {Consulter le journal d'audit};

  % --- associations principales
  \foreach \a in {dir,resp,contrib,adm,audit}
    \draw[lien] (\a.east) -- (connect.west);
  \foreach \a in {dsi,data,ai,fin,pmo}
    \draw[lien] (\a.west) -- (connect.east);

  \draw[lien] (dir.east)   -- (prio.west);
  \draw[lien] (dir.east)   -- (valid.west);
  \draw[lien] (dir.east)   -- (roi.west);
  \draw[lien] (resp.east)  -- (opp.west);
  \draw[lien] (resp.east)  -- (diag.west);
  \draw[lien] (contrib.east) -- (repond.west);
  \draw[lien] (contrib.east) -- (opp.west);
  \draw[lien] (adm.east)   -- (admin.west);
  \draw[lien] (audit.east) -- (journal.west);

  \draw[lien] (dsi.west)   -- (qualif.east);
  \draw[lien] (data.west)  -- (repond.east);
  \draw[lien] (ai.west)    -- (over.east);
  \draw[lien] (ai.west)    -- (reco.east);
  \draw[lien] (fin.west)   -- (roi.east);
  \draw[lien] (pmo.west)   -- (suivi.east);
  \draw[lien] (pmo.west)   -- (fiche.east);

  % --- dépendance interne
  \draw[-{Stealth[length=4pt]}, dashed, draw=medicaGreen] (opp) -- node[midway, above, font=\tiny, text=medicaGreen] {\guillemotleft include\guillemotright} (qualif);
\end{tikzpicture}
\caption{Diagramme de cas d'utilisation général}
\label{fig:uc-general}
\end{figure}

La relation \guillemotleft~include~\guillemotright{} entre la création d'une opportunité
et sa qualification n'est pas décorative : dans le code, l'endpoint de création appelle
systématiquement le moteur de qualification avant de répondre. Il n'existe pas
d'opportunité non qualifiée.

\subsection{Sous-système Diagnostic et Readiness}

\begin{figure}[H]
\centering
\begin{tikzpicture}[scale=0.95, transform shape]
  \node[acteur] (resp) at (0,2.2) {Responsable de département};
  \node[acteur] (dsi)  at (0,0.8) {DSI / Data Office};
  \node[acteur] (dir)  at (0,-0.6) {Direction};

  \begin{scope}[on background layer]
    \draw[draw=medicaBorder, fill=medicaBg, rounded corners=8pt] (2.8,-1.6) rectangle (10.4,3.1);
  \end{scope}
  \node[font=\small\bfseries, text=medicaDeep] at (6.6,2.75) {Module M2 / M5};

  \node[cas] (camp)  at (4.6,2.0) {Ouvrir une campagne};
  \node[cas] (note)  at (4.6,1.0) {Noter une dimension};
  \node[cas] (preuve)at (4.6,0.0) {Joindre une preuve};
  \node[cas] (calc)  at (8.4,1.5) {Calculer le readiness};
  \node[cas] (bloq)  at (8.4,0.5) {Identifier les dimensions bloquantes};
  \node[cas] (exp)   at (8.4,-0.6) {Exporter le diagnostic};
  \node[cas] (valcamp) at (4.6,-1.0) {Valider la campagne};

  \draw[lien] (resp.east) -- (camp.west);
  \draw[lien] (resp.east) -- (valcamp.west);
  \draw[lien] (dsi.east)  -- (note.west);
  \draw[lien] (dsi.east)  -- (preuve.west);
  \draw[lien] (dir.east)  -- (exp.west);
  \draw[lien] (dir.east)  -- (calc.west);

  \draw[-{Stealth[length=4pt]}, dashed, draw=medicaGreen] (note) -- node[midway, above, font=\tiny, text=medicaGreen] {\guillemotleft include\guillemotright} (calc);
  \draw[-{Stealth[length=4pt]}, dashed, draw=medicaGreen] (calc) -- node[midway, right, font=\tiny, text=medicaGreen] {\guillemotleft extend\guillemotright} (bloq);
\end{tikzpicture}
\caption{Cas d'utilisation du diagnostic et du calcul de readiness}
\label{fig:uc-diagnostic}
\end{figure}

\subsection{Sous-système Portefeuille et Valeur}

\begin{figure}[H]
\centering
\begin{tikzpicture}[scale=0.95, transform shape]
  \node[acteur] (dir) at (0,2.2) {Direction};
  \node[acteur] (fin) at (0,0.8) {Finance};
  \node[acteur] (pmo) at (0,-0.6) {PMO};

  \begin{scope}[on background layer]
    \draw[draw=medicaBorder, fill=medicaBg, rounded corners=8pt] (2.8,-2.2) rectangle (10.4,3.1);
  \end{scope}
  \node[font=\small\bfseries, text=medicaDeep] at (6.6,2.75) {Modules M7 / M8 / M9 / M10};

  \node[cas] (prio)  at (4.6,2.0) {Ajuster les pondérations};
  \node[cas] (recalc)at (8.4,2.0) {Recalculer le classement};
  \node[cas] (fiche) at (4.6,1.0) {Générer la fiche projet};
  \node[cas] (dec)   at (8.4,1.0) {Rendre la décision de comité};
  \node[cas] (cout)  at (4.6,0.0) {Saisir coûts et gains};
  \node[cas] (bc)    at (8.4,0.0) {Calculer le business case};
  \node[cas] (gate)  at (4.6,-1.0) {Passer une gate};
  \node[cas] (suivi) at (8.4,-1.0) {Suivre l'avancement};
  \node[cas] (roi)   at (6.5,-1.9) {Comparer valeur promise et réalisée};

  \draw[lien] (dir.east) -- (prio.west);
  \draw[lien] (dir.east) -- (fiche.west);
  \draw[lien] (fin.east) -- (cout.west);
  \draw[lien] (fin.east) -- (roi.west);
  \draw[lien] (pmo.east) -- (gate.west);
  \draw[lien] (pmo.east) -- (suivi.west);

  \draw[-{Stealth[length=4pt]}, dashed, draw=medicaGreen] (prio) -- node[midway, above, font=\tiny, text=medicaGreen] {\guillemotleft include\guillemotright} (recalc);
  \draw[-{Stealth[length=4pt]}, dashed, draw=medicaGreen] (cout) -- node[midway, above, font=\tiny, text=medicaGreen] {\guillemotleft include\guillemotright} (bc);
  \draw[-{Stealth[length=4pt]}, dashed, draw=medicaGreen] (fiche) -- node[midway, above, font=\tiny, text=medicaGreen] {\guillemotleft extend\guillemotright} (dec);
\end{tikzpicture}
\caption{Cas d'utilisation de la priorisation, de la fiche projet et de la valeur}
\label{fig:uc-portefeuille}
\end{figure}

\section{Modèle de données}

Le schéma compte \textbf{34~tables}, dont quatre tables d'association. Les diagrammes qui
suivent le découpent en quatre domaines pour rester lisibles. Les noms sont ceux des
classes SQLAlchemy réelles, la table physique étant précisée lorsqu'elle diffère.

\subsection{Domaine A — Organisation et contrôle d'accès}

\begin{figure}[H]
\centering
\begin{tikzpicture}[scale=0.9, transform shape, every node/.style={font=\scriptsize}]
  \node[classe] (org) at (0,0) {\umlclasse{Organization}{
    \underline{id} : int\\ name, code (unique)\\ industry, country\\ description}};
  \node[classe] (site) at (-4.6,-2.9) {\umlclasse{Site}{
    \underline{id} : int\\ \textit{organization\_id} FK\\ name, code\\ city, country, site\_type}};
  \node[classe] (dept) at (0,-2.9) {\umlclasse{Department}{
    \underline{id} : int\\ \textit{organization\_id} FK\\ \textit{site\_id} FK\\ name, code, family\\ head\_name}};
  \node[classe] (user) at (5,-2.9) {\umlclasse{User (app\_user)}{
    \underline{id} : int\\ \textit{organization\_id} FK\\ \textit{department\_id} FK\\ email (unique)\\ full\_name, job\_title\\ hashed\_password\\ is\_active, is\_superuser\\ last\_login\_at}};
  \node[classe] (role) at (5,-6.6) {\umlclasse{Role}{
    \underline{id} : int\\ \textit{organization\_id} FK\\ code, label\\ scope\_level\\ is\_system, is\_platform}};
  \node[classe] (perm) at (0,-6.6) {\umlclasse{Permission}{
    \underline{id} : int\\ code (unique)\\ label, module, group}};

  \draw[lien] (org.west) -| node[near start, above, font=\tiny] {1..*} (site.north);
  \draw[lien] (org.south) -- node[midway, right, font=\tiny] {1..*} (dept.north);
  \draw[lien] (org.east) -| node[near start, above, font=\tiny] {1..*} (user.north);
  \draw[lien] (org.east) -| ([xshift=0.6cm]role.north);
  \draw[lien] (dept.east) -- node[midway, above, font=\tiny] {0..1} (user.west);
  \draw[lien] (user.south) -- node[midway, right, font=\tiny] {\textit{user\_role} *..*} (role.north);
  \draw[lien] (role.west) -- node[midway, above, font=\tiny] {\textit{role\_permission} *..*} (perm.east);
  \draw[lien, dashed] (user.west) -- node[midway, below, font=\tiny] {\textit{user\_department} *..*} (dept.east);
\end{tikzpicture}
\caption{Domaine A — Organisation et contrôle d'accès}
\label{fig:classes-rbac}
\end{figure}

Quatre points méritent commentaire.

\begin{itemize}
  \item La table \code{User} s'appelle physiquement \code{app\_user}, \code{user} étant
    un mot réservé en PostgreSQL.
  \item Les permissions sont \emph{globales à la plateforme} : seuls les rôles, qui les
    regroupent, portent un \code{organization\_id}. Deux organisations partagent donc le
    même catalogue de permissions mais ont chacune leurs rôles.
  \item \code{User.department\_id} coexiste avec la table d'association
    \code{user\_department}. Le premier désigne le département principal et sert de
    repli ; la seconde porte les rattachements multiples. Cette coexistence est le
    résultat d'une migration qui devait préserver les données existantes.
  \item \code{Role.scope\_level} porte le périmètre par défaut. Le périmètre effectif
    d'un utilisateur est le plus large de ceux de ses rôles.
\end{itemize}

\subsection{Domaine B — Diagnostic, architecture et opportunités}

\begin{figure}[H]
\centering
\begin{tikzpicture}[scale=0.88, transform shape, every node/.style={font=\scriptsize}]
  \node[classe] (ass) at (0,0) {\umlclasse{Assessment}{
    \underline{id} : int\\ \textit{organization\_id} FK\\ \textit{site\_id}, \textit{department\_id} FK\\ name, campaign, status\\ blocking\_threshold\\ completion\_rate}};
  \node[classe] (resp) at (0,-3.4) {\umlclasse{AssessmentResponse}{
    \underline{id} : int\\ \textit{assessment\_id} FK\\ axis, dimension\_code\\ dimension\_label, question\\ score (0--5), weight\\ confidence (0--1)\\ evidence, is\_critical}};
  \node[classe] (opp) at (5.8,-1.2) {\umlclasse{Opportunity}{
    \underline{id} : int\\ \textit{organization\_id} FK\\ \textit{department\_id}, \textit{site\_id} FK\\ \textit{process\_id} FK\\ code, title, summary\\ as\_is\_context, pain\_points\\ answers : JSON\\ 15 faits de qualification\\ 6 composantes de score\\ risk\_probability, risk\_impact\\ priority\_score, rank, wave\\ score\_explanation : JSON}};
  \node[classe] (reco) at (11.4,0.2) {\umlclasse{AIRecommendation}{
    \underline{id} : int\\ \textit{opportunity\_id} FK (1:1)\\ ai\_type, ai\_type\_label\\ solution\_pattern\\ rationale, confidence\\ fired\_rules : JSON\\ alternatives : JSON\\ guardrails : JSON}};
  \node[classe] (dep) at (11.4,-3.4) {\umlclasse{Dependency}{
    \underline{id} : int\\ \textit{source\_opportunity\_id} FK\\ \textit{target\_opportunity\_id} FK\\ kind, asset\_ref, label\\ strength, is\_blocking}};
  \node[classe] (arch) at (5.8,-5.9) {\umlclasse{Process / Application / DataSource / KPI}{
    Référentiel AS-IS, chacun rattaché à l'organisation\\
    DataSource : quality\_score, has\_ground\_truth\\
    KPI : baseline, target, current}};

  \draw[lien] (ass.south) -- node[midway, right, font=\tiny] {1..*} (resp.north);
  \draw[lien] (opp.east) -- node[midway, above, font=\tiny] {1..1} (reco.west);
  \draw[lien] (opp.east) -| node[near start, below, font=\tiny] {*..*} (dep.west);
  \draw[lien] (opp.south) -- node[midway, right, font=\tiny] {0..1} (arch.north);
\end{tikzpicture}
\caption{Domaine B — Diagnostic, référentiel d'architecture et opportunités}
\label{fig:classes-diagnostic}
\end{figure}

La table \code{opportunity} est la plus large du schéma avec 48~colonnes. Elle regroupe
trois familles d'information : le contexte métier rédigé, quinze faits booléens ou
numériques qui alimentent le moteur de qualification, et six composantes de score
assorties du résultat calculé. Le champ \code{score\_explanation}, de type JSON,
conserve la décomposition du score au moment du dernier calcul.

\code{AIRecommendation} entretient une relation un à un avec l'opportunité. Son nom
contient « AI » parce qu'elle décrit un type de solution d'intelligence artificielle ;
elle est produite par un moteur de règles, non par un modèle.

\subsection{Domaine C — Projets et valeur}

\begin{figure}[H]
\centering
\begin{tikzpicture}[scale=0.88, transform shape, every node/.style={font=\scriptsize}]
  \node[classe] (proj) at (0,0) {\umlclasse{Project}{
    \underline{id} : int\\ \textit{organization\_id} FK\\ \textit{opportunity\_id} FK (1:1)\\ \textit{department\_id} FK\\ code, name, phase, status\\ sponsor, owners\\ budget\_eur, committed\_eur\\ progress, next\_gate\_date}};
  \node[classe] (stage) at (5.4,1.4) {\umlclasse{Stage}{
    \underline{id} : int\\ \textit{project\_id} FK\\ code, name, sequence\\ status, gate\_passed\\ gate\_criteria : JSON}};
  \node[classe] (ms) at (5.4,-1.0) {\umlclasse{Milestone}{
    \underline{id} : int\\ \textit{project\_id}, \textit{stage\_id} FK\\ name, due\_date\\ status, is\_gate}};
  \node[classe] (task) at (5.4,-3.2) {\umlclasse{Task}{
    \underline{id} : int\\ \textit{project\_id}, \textit{stage\_id} FK\\ title, assignee\\ status, due\_date, effort}};
  \node[classe] (risk) at (0,-3.6) {\umlclasse{Risk}{
    \underline{id} : int\\ \textit{project\_id} FK\\ \textit{opportunity\_id} FK\\ title, category\\ probability (1--5)\\ impact (1--5)\\ severity, status}};
  \node[classe] (mit) at (0,-6.4) {\umlclasse{Mitigation}{
    \underline{id} : int\\ \textit{risk\_id} FK\\ action, owner\\ due\_date, status\\ residual\_prob/impact}};
  \node[classe] (cost) at (10.6,0.6) {\umlclasse{Cost}{
    \underline{id} : int\\ \textit{project\_id} FK\\ \textit{opportunity\_id} FK\\ label, category, nature\\ amount\_eur, recurring}};
  \node[classe] (ben) at (10.6,-2.2) {\umlclasse{Benefit}{
    \underline{id} : int\\ \textit{opportunity\_id} FK\\ \textit{kpi\_id} FK\\ annual\_amount\_eur\\ realization\_rate\\ pnl\_line}};
  \node[classe] (meas) at (10.6,-5.0) {\umlclasse{Measurement}{
    \underline{id} : int\\ \textit{kpi\_id} FK\\ measured\_on, value\\ period\_label, source}};

  \draw[lien] (proj.east) -- node[midway, above, font=\tiny] {1..*} (stage.west);
  \draw[lien] (proj.east) -- node[midway, above, font=\tiny] {1..*} (ms.west);
  \draw[lien] (proj.east) -- node[midway, below, font=\tiny] {1..*} (task.west);
  \draw[lien] (proj.south) -- node[midway, right, font=\tiny] {1..*} (risk.north);
  \draw[lien] (risk.south) -- node[midway, right, font=\tiny] {1..*} (mit.north);
  \draw[lien] (stage.east) -| (cost.north);
  \draw[lien] (ben.west) -- node[midway, above, font=\tiny] {0..1} (task.east);
\end{tikzpicture}
\caption{Domaine C — Cycle de vie des projets et valorisation}
\label{fig:classes-projets}
\end{figure}

\code{Cost} et \code{Benefit} peuvent être rattachés soit à une opportunité, soit à un
projet. Dans le jeu de démonstration ils le sont à l'opportunité, ce qui permet de
calculer un business case avant même que le projet n'existe.

\subsection{Domaine D — Gouvernance documentaire et audit}

\begin{figure}[H]
\centering
\begin{tikzpicture}[scale=0.92, transform shape, every node/.style={font=\scriptsize}]
  \node[classe] (doc) at (0,0) {\umlclasse{Document}{
    \underline{id} : int\\ \textit{organization\_id} FK\\ \textit{opportunity\_id} FK\\ \textit{project\_id} FK\\ code, title, doc\_type\\ status, current\_version\\ author}};
  \node[classe] (ver) at (5.4,1.2) {\umlclasse{Version}{
    \underline{id} : int\\ \textit{document\_id} FK\\ number, label\\ content : JSON\\ change\_summary\\ created\_by, created\_at}};
  \node[classe] (app) at (5.4,-1.6) {\umlclasse{Approval}{
    \underline{id} : int\\ \textit{document\_id} FK\\ \textit{opportunity\_id} FK\\ \textit{project\_id} FK\\ \textit{approver\_user\_id} FK\\ stage\_code, committee\\ decision, decided\_on}};
  \node[classe] (log) at (0,-4.2) {\umlclasse{AuditLog — insertion seule}{
    \underline{id} : int\\ organization\_id, user\_id\\ user\_email\\ action, entity\_type\\ entity\_id, entity\_label\\ old\_value : JSON\\ new\_value : JSON\\ ip\_address, created\_at}};
  \node[classe] (notif) at (6,-4.6) {\umlclasse{Notification}{
    \underline{id} : int\\ \textit{organization\_id} FK\\ \textit{user\_id} FK\\ title, body, level\\ entity\_type, entity\_id\\ is\_read}};

  \draw[lien] (doc.east) -- node[midway, above, font=\tiny] {1..*} (ver.west);
  \draw[lien] (doc.east) -- node[midway, below, font=\tiny] {1..*} (app.west);
  \draw[lien, dashed, draw=medicaAmber] (log.north) -- node[midway, right, font=\tiny, text=medicaAmber]
    {trace toute écriture} (doc.south);
\end{tikzpicture}
\caption{Domaine D — Documents versionnés, approbations et journal d'audit}
\label{fig:classes-audit}
\end{figure}

\code{AuditLog} est volontairement dépourvue de clés étrangères. Une entrée d'audit doit
survivre à la suppression de l'entité qu'elle décrit : si l'opportunité est effacée, la
trace de sa suppression doit subsister. Les identifiants y sont donc de simples entiers,
sans contrainte référentielle.

\section{Architecture technique}

\subsection{Vue d'ensemble}

\begin{figure}[H]
\centering
\begin{tikzpicture}[node distance=0.65cm, scale=0.95, transform shape]
  \node[blocsombre, text width=3.4cm] (user) {\textbf{Utilisateur}\\\scriptsize navigateur};

  \node[bloc, below=of user, text width=10.5cm, fill=medicaSoft, draw=medicaGreen] (front)
    {\textbf{Frontend React 18 + TypeScript} \\[2pt]
     \scriptsize Vite · React Router · TanStack Query · Recharts · axios};

  \node[bloc, below=of front, text width=10.5cm, fill=white, draw=medicaTeal] (api)
    {\textbf{API REST — 99 endpoints sous \code{/api/v1}} \\[2pt]
     \scriptsize JSON · OpenAPI sur \code{/docs} · CORS restreint aux origines déclarées};

  \node[bloc, below=of api, text width=10.5cm, fill=white, draw=medicaGreen] (fastapi)
    {\textbf{FastAPI} — routers par module (M1 à M12) \\[2pt]
     \scriptsize Dépendances : session, utilisateur courant, permission, périmètre};

  \node[bloc, below=of fastapi, text width=10.5cm, fill=medicaSoft, draw=medicaGreen] (services)
    {\textbf{Services métier} — sans dépendance à FastAPI \\[2pt]
     \scriptsize readiness · recommendation · scoring · crossing · value · questionnaire · audit};

  \node[bloc, below=of services, text width=10.5cm, fill=white, draw=medicaTeal] (orm)
    {\textbf{SQLAlchemy 2.0} — 34 tables, sessions par requête \\[2pt]
     \scriptsize Migrations versionnées par Alembic};

  \node[blocsombre, below=of orm, text width=10.5cm] (db)
    {\textbf{PostgreSQL} \\[2pt]\scriptsize Filtrage \code{organization\_id} sur toute table métier};

  \draw[flechevert] (user) -- (front);
  \draw[flechevert] (front) -- node[right, font=\scriptsize, text=medicaMuted] {HTTPS + jeton JWT} (api);
  \draw[flechevert] (api) -- (fastapi);
  \draw[flechevert] (fastapi) -- (services);
  \draw[flechevert] (services) -- (orm);
  \draw[flechevert] (orm) -- (db);

  % éléments transverses
  \node[bloc, right=1.1cm of fastapi, text width=2.6cm, fill=medicaBg, draw=medicaAmber, font=\scriptsize]
    (sec) {\textbf{Sécurité}\\ JWT HS256\\ bcrypt\\ RBAC\\ périmètre};
  \draw[fleche, dashed, draw=medicaAmber] (sec) -- (fastapi);

  \node[bloc, right=1.1cm of db, text width=2.6cm, fill=medicaBg, draw=medicaTeal, font=\scriptsize]
    (docker) {\textbf{Docker}\\ postgres:16\\ backend\\ frontend};
  \draw[fleche, dashed, draw=medicaTeal] (docker) -- (db);
\end{tikzpicture}
\caption{Architecture en couches de la plateforme}
\label{fig:architecture}
\end{figure}

\subsection{Frontend}

L'interface est une application monopage construite avec React~18 et TypeScript,
servie par Vite. Le routage est assuré par React Router, les appels réseau par axios
encapsulé dans des hooks TanStack Query, et les graphiques par Recharts.

Le point structurant est que \emph{la navigation dérive des permissions}. Le composant
\code{Layout} filtre les entrées de menu selon les permissions renvoyées par
\code{/auth/me} ; une entrée n'apparaît que si son écran est réellement accessible. Le
même principe gouverne les boutons d'action.

\subsection{Backend}

L'API est écrite en Python~3.12 avec FastAPI. Elle est organisée en dix routers, un par
domaine fonctionnel, montés sous le préfixe \code{/api/v1}.

La logique métier vit dans sept services qui n'importent rien de FastAPI. Cette
séparation a une conséquence pratique : les formules de readiness, de score et de ROI
se testent en appelant directement les fonctions, sans serveur HTTP ni base de données.

\subsection{Base de données}

PostgreSQL héberge les 34~tables. Le schéma est versionné par trois migrations Alembic,
écrites pour s'appliquer sur une base déjà peuplée. La troisième, qui a introduit les
périmètres, crée les colonnes avec une valeur par défaut serveur le temps de la reprise,
recopie les rattachements existants depuis \code{department\_id}, puis retire la valeur
par défaut.

\subsection{API}

L'API compte 99~endpoints répartis comme suit.

\begin{table}[H]
\centering
\caption{Répartition des endpoints par domaine}
\label{tab:endpoints}
\small
\begin{tabular}{@{}L{3.6cm} C{1.6cm} L{8.2cm}@{}}
\toprule
\textbf{Préfixe} & \textbf{Nombre} & \textbf{Objet} \\
\midrule
\code{/org}           & 18 & Organisation, sites, départements, utilisateurs, rôles \\
\code{/projects}      & 16 & Projets, phases, jalons, tâches, risques, gates \\
\code{/crossings}     & 14 & Processus, applications, données, dépendances, doublons \\
\code{/opportunities} & 12 & Opportunités, qualification, priorisation \\
\code{/value}         &  9 & Coûts, bénéfices, mesures, business case \\
\code{/diagnostic}    &  8 & Campagnes, réponses, readiness \\
\code{/sheets}        &  7 & Fiches projet, versions, approbations, export \\
\code{/auth}          &  5 & Connexion, profil, mot de passe \\
\code{/admin}         &  4 & Journal d'audit, notifications \\
\code{/dashboards}    &  4 & Agrégations Direction, département, portefeuille \\
Système               &  2 & Sonde de santé et racine \\
\midrule
\textbf{Total}        & \textbf{99} & \\
\bottomrule
\end{tabular}
\end{table}

\subsection{Déploiement}

Un fichier \code{docker-compose.yml} décrit trois services : PostgreSQL~16 avec sonde
de santé et volume persistant, le backend construit depuis \code{backend/Dockerfile}, et
le frontend depuis \code{frontend/Dockerfile}. Le backend attend que la base soit saine
avant de démarrer.

\begin{remarque}[Sur la version de PostgreSQL]
Le fichier Docker déclare \code{postgres:16-alpine}. L'instance utilisée pendant le
développement est une PostgreSQL~12.4 installée localement. Aucune fonctionnalité
spécifique à la version~16 n'est employée, mais cet écart mérite d'être connu avant une
mise en production.
\end{remarque}
